%==============================================================================
% IEEE 期刊论文：HARP-FDN-Optimized
% 编译: xelatex main_ieee.tex（推荐，支持中文）
%       xelatex main_ieee && xelatex main_ieee（两遍以解析交叉引用）
% ==============================================================================
\documentclass[journal]{IEEEtran}

% ---- 中文支持 ----
\usepackage{xeCJK}
\setCJKmainfont{SimSun}[
  BoldFont=SimHei,
  ItalicFont=KaiTi,
]
\setCJKsansfont{SimHei}
\setCJKmonofont{FangSong}

% ---- 数学 ----
\usepackage{amsmath,amssymb,bm}

% ---- 图形 ----
\usepackage{graphicx}
\usepackage{svg}
\graphicspath{{figures/}}

% ---- 表格 ----
\usepackage{booktabs,multirow,array}
\usepackage{stfloats}

% ---- URL ----
\usepackage{url}

% ---- 引用 ----
\usepackage{cite}

% ---- 列表 ----
\usepackage{enumitem}
\setlist[itemize]{noitemsep,topsep=2pt,leftmargin=1.25em}
\setlist[enumerate]{noitemsep,topsep=2pt,leftmargin=1.25em}

\hyphenation{op-tical net-works semi-conduc-tor}

\begin{document}

% ==============================================================================
% 标题与作者
% ==============================================================================
\title{HARP-FDN-Optimized：基于共享概率混合骨干与后验校准的原型预报分解网络}

\author{作者A,~\IEEEmembership{学生会员,~IEEE,}
        作者B,~\IEEEmembership{会员,~IEEE,}
        and~作者C,~\IEEEmembership{高级会员,~IEEE}}
\thanks{作者A 就读于 XX 大学计算机学院，XX 市，XX 国（电子邮箱：authora@university.edu）。}
\thanks{作者B 和作者C 就职于 XX 大学电子工程学院，XX 市，XX 国（电子邮箱：\{authorb, authorc\}@university.edu）。}
\thanks{收稿日期：XXXX年XX月；修订日期：XXXX年XX月。}

\markboth{IEEE Transactions on Neural Networks and Learning Systems, Vol.~XX, No.~XX, XXXX}%
{作者A \MakeLowercase{\textit{等}}：HARP-FDN-Optimized}

% ==============================================================================
% 摘要
% ==============================================================================
\begin{abstract}
基于非均匀支撑点的预报分解网络已经在时间序列与光谱回归任务中展现出较好的稳定性，但原方法在长步预测、跨尺度不确定性利用以及编码结构效率方面仍存在三点不足：第一，细支撑解码直接依赖绝对支撑点，容易在长预测区间产生漂移；第二，细粒度与粗粒度分支各自维护独立编码器，参数与计算冗余显著；第三，支撑点后验中的方差信息主要被用作分类置信度，尚未充分参与融合与校准。本文提出 HARP-FDN-Optimized，面向同一骨干输入同时输出 fine 双分支与 coarse 双分支支撑对数，并在后端统一执行相对原型解码、不确定性感知融合与后验校准。具体而言，本文在原有框架上做出三项结构化改进：其一，引入共享 PatchMixer 骨干，以短期 patch 流与长期 patch 流的尺度门控融合替代四组独立编码器，在保持概率混合建模能力的同时大幅降低参数量；其二，在细分支上显式引入相对原型解码，使 fine 输出由 coarse 原型锚点和局部残差修正共同确定，并通过锚点方差自适应调节相对解码强度；其三，将跨尺度一致性约束从点值 MSE 扩展为均值-方差分布对齐，同时在后验不确定区域引入弱残差校准，从而更充分地利用支撑后验中的分布信息。HARP-FDN-Optimized 保留了原有后验概率建模范式，但将四个独立预测网络压缩为一个共享骨干与四个轻量后验头，并在解码端统一执行 coarse-to-fine 的相对解码、不确定性融合以及后验校准，使模型在保持可解释性的同时获得更稳定的长步预测与更紧凑的结构。
\end{abstract}

\begin{IEEEkeywords}
预报分解网络，共享 PatchMixer 骨干，相对原型解码，分布一致性约束，后验校准，时间序列预测，CEMP 恒星参数回归。
\end{IEEEkeywords}

\IEEEpeerreviewmaketitle

% ==============================================================================
% 一、引言
% ==============================================================================
\section{引言}
\label{sec:introduction}

序列数据广泛存在于工程观测与科学分析中。电力负荷、交通流量、设备状态与工业过程时间序列需要模型同时捕获长期趋势、短期波动与非平稳变化；天文光谱序列则承载恒星大气的物理与化学信息，通过反演有效温度 $T_{\mathrm{eff}}$、表面重力 $\log g$ 与金属丰度 $[\mathrm{Fe/H}]$ 等参数支撑恒星演化研究，其中 CEMP（Carbon-Enhanced Metal-Poor）恒星参数回归是代表性的高维光谱建模问题。上述任务虽然目标形式不同，但本质上都属于基于序列观测的多变量建模问题：模型需要在长依赖、高噪声以及分布偏移条件下保持预测稳定性与不确定性建模能力。

近年来，深度学习在时间序列建模上取得了显著进展。从 RNN/LSTM、TCN 到 MLP 与 Transformer 系列方法，预测精度持续提升；Patch 分块~\cite{patchtst}、可逆归一化~\cite{revin} 与跨变量 token 化~\cite{itransformer} 等技术进一步增强了模型对非平稳序列的适应能力。然而，在长步预测任务中，现有基于支撑点的预报网络仍面临两个核心问题：第一，绝对解码方式在预测 horizon 增长时容易累积偏移，尤其当细支撑点分布较宽或趋势变化剧烈时更为明显；第二，支撑后验中的方差与置信度信息尚未被充分纳入融合与校准机制，导致不确定性估计与最终输出之间存在信息损失。

针对上述问题，本文在 HARP-FDN 的基础上提出 HARP-FDN-Optimized。该方法保留原始框架的 RevIN 归一化~\cite{revin}、趋势-季节分解以及非均匀支撑点概率预测范式，但在结构上做出三点关键改进。首先，模型将原本四个独立 PatchMixer 编码器替换为一个共享概率混合骨干，以短期 patch 流、长期 patch 流和输入条件尺度门控构建统一特征，并通过四个轻量后验头分别输出 fine 双分支与 coarse 双分支支撑对数，在降低参数冗余的同时保留多尺度建模能力。其次，模型在细分支上显式引入相对原型解码，以 coarse 融合输出为锚点，结合锚点方差与可学习相对强度参数动态调节 fine 分支对局部残差修正的依赖程度，从而在长 horizon 上获得更稳定的解码轨迹。最后，模型将跨尺度约束从点值一致性扩展为均值-方差分布对齐，并引入后验门控的残差校准，使输出在高不确定区域获得保守修正，在低不确定区域保持主预测结构不变。

本文的主要贡献概括如下：
\begin{itemize}
  \item 提出共享概率混合骨干结构，以 short/long PatchMixer 流和输入条件尺度门控替代四组独立编码器，在保持支撑对数生成能力的同时显著降低模型复杂度。
  \item 设计相对原型解码机制，以 coarse 融合输出为未来轨迹锚点，fine 分支预测相对于该锚点的残差分布，并通过锚点方差与可学习门控自适应调节相对解码强度。
  \item 将跨尺度一致性约束拓展为分布对齐，使 fine 与 coarse 预测在均值与方差上保持一致，从而更充分地利用支撑后验中的不确定性信息。
  \item 引入后验门控残差校准，仅在模型高不确定区域对输出施加弱修正，在不破坏主预测结构的前提下缓解离散支撑点带来的量化偏差。
\end{itemize}

% ==============================================================================
% 二、相关工作
% ==============================================================================
\section{相关工作}
\label{sec:related_work}

\subsection{长期时间序列预测}
长期预测方法主要围绕 Transformer 结构优化、分解先验与输入表示展开。Informer~\cite{informer} 通过稀疏注意力降低长序列计算复杂度，Autoformer~\cite{autoformer} 与 FEDformer~\cite{fedformer} 分别引入自相关与频域分解以增强周期建模能力，iTransformer~\cite{itransformer} 将变量维 token 化，使 Transformer Encoder 能够直接建模跨变量依赖。这些方法在不同层面提升了长步预测能力，但大多基于点值输出或隐式统计建模，难以提供显式的概率支撑与不确定性表达。

\subsection{支撑点与概率预报}
支撑点方法将连续预测问题转化为在有限支撑集上的概率估计问题，具有建模简洁、可解释性较强的优势。本文沿用的预报分解网络进一步将该思路与趋势-季节分解结合，使模型在 coarse 与 fine 两个尺度上同时建模未来轨迹。然而，原始框架在长步解码时仍以绝对支撑期望为主，且 coarse 与 fine 之间的约束主要停留在点值层面，导致后验中蕴含的方差信息未被充分使用。

\subsection{条件调制与归一化机制}
FiLM~\cite{film} 等条件调制方法通过可学习的缩放与平移控制特征分布，在图像、语言和时序任务中均被广泛使用。RevIN~\cite{revin} 通过可逆归一化处理分布漂移，已成为非平稳时间序列建模的重要预处理手段。本文在此基础上，将相对解码、分布一致性与后验校准统一纳入 HARP-FDN 框架，使模型在保留可逆归一化与分解先验的同时，进一步提升对不确定性与长步漂移的控制能力。

% ==============================================================================
% 三、方法
% ==============================================================================
\section{方法}
\label{sec:method}

本节以 HARP-FDN-Optimized 的实际实现为依据描述模型结构。该模型延续了 RevIN 归一化~\cite{revin}、趋势-季节分解以及双支撑集概率预测的整体范式，但将核心编码结构由四组独立 PatchMixer 压缩为一个共享骨干，并在解码端引入相对原型解码、不确定性感知融合与后验校准。

\subsection{问题形式化}
给定历史序列 $\mathbf{X} \in \mathbb{R}^{B \times L \times C}$，模型需要输出未来预测 $\hat{\mathbf{Y}} \in \mathbb{R}^{B \times H \times C}$，其中 $B$、$L$、$H$、$C$ 分别表示批量大小、历史长度、预测长度和变量数。HARP-FDN-Optimized 在内部进一步维护两组非均匀支撑集 $\mathbf{S}_1, \mathbf{S}_2 \in \mathbb{R}^{N_s \times 1}$，并在 coarse 尺度上使用长度为 $L_c$ 的短期原型预测，再通过上采样与残差修正得到最终细尺度输出。

\subsection{归一化与分解}
模型首先使用 RevIN~\cite{revin} 对输入进行可逆实例归一化，以缓解序列间的分布偏移。随后，根据移动平均类型将归一化结果分解为季节分量与趋势分量。该分解不是全局固定的，而是随输入条件自适应调整，从而使后续骨干网络能够在更清晰的结构特征上进行学习。

\subsection{共享概率混合骨干}
\label{subsec:backbone}
原 HARP-FDN 使用四个独立 PatchMixer 编码器分别生成 fine 双分支与 coarse 双分支支撑对数。HARP-FDN-Optimized 将其替换为一个共享概率混合骨干，并引入四个轻量后验头。骨干结构由以下部分组成：
\begin{enumerate}[leftmargin=*,itemsep=2pt]
  \item \textbf{短期 PatchMixer 流}：采用较小 patch 长度与步长，用于刻画局部波动、短期变形和高频结构。
  \item \textbf{长期 PatchMixer 流}：采用标准 patch 长度与步长，用于刻画周期片段、长期趋势与跨段依赖。
  \item \textbf{输入条件尺度门控}：基于两流 summary 特征生成尺度权重，对短期与长期 horizon 预测进行自适应加权融合。
  \item \textbf{季节线性路径与趋势门控路径}：分别保留显式季节线性投影和趋势门控外推，使骨干在 patch 混合建模之外仍保持对季节趋势结构的直接访问能力。
\end{enumerate}
骨干输出被归一化为统一特征，并分别送入四个后验头，对应 fine 与 coarse 两类尺度、以及两个支撑集维度。该设计在参数上实现了显著压缩，同时保留了对支撑分布与多尺度行为的表达能力。

\subsection{后验头与支撑对数生成}
每个后验头均由 LayerNorm、两层 MLP 与支撑维度重塑操作组成，输出形状为 $B \times H \times N_s \times C$ 的支撑对数。与独立编码器方案相比，该结构使 fine 与 coarse 分支共享骨干表征，仅在最后阶段根据支撑集与尺度需求进行分化，从而在保持预测多样性的同时避免重复建模。

\subsection{相对原型解码}
\label{subsec:relative}
支撑对数经过温度缩放与 softmax 后，可得到支撑点上的概率分布，并通过期望解码得到绝对预测值与预测方差。HARP-FDN-Optimized 在此基础上引入相对原型解码：首先将 coarse 双分支融合结果视为未来轨迹原型，随后在 fine 分支中预测相对于该原型的残差支撑分布，并通过锚点方差和可学习相对强度参数控制模型在绝对解码与相对修正之间的切换比例。该机制使模型在长步预测中先建立整体轨迹锚点，再进行局部修正，从而减轻直接绝对解码带来的漂移。

\subsection{不确定性感知融合与分布一致性}
对于同尺度的双分支输出，模型首先根据支撑熵、最大概率与方差估计分支置信度，并据此生成软融合权重。更关键的是，HARP-FDN-Optimized 将跨尺度约束从点值 MSE 推广为分布一致性：先将 fine 预测均值与方差聚合到 coarse 时间尺度，再要求聚合后的 fine 分布与 coarse 分布在均值与方差上同时匹配。该设计能够更充分地利用支撑后验中的不确定性信息，使 coarse 原型对 fine 输出形成更完整的结构性约束。

\subsection{后验门控残差校准}
\label{subsec:calibration}
离散支撑点解码可能引入量化偏差。为此，模型增加一个基于历史输入的弱残差校准分支，并使用后验不确定性对其进行门控。具体做法是先根据预测方差生成不确定性门，再将其与可学习校准强度相乘，最终以较小幅度修正主输出。该设计保证模型在低不确定区域几乎不改变主预测结构，而在高不确定区域则引入保守修正，从而在稳定性与精度之间取得更好的平衡。

\subsection{训练目标}
训练过程同时优化主预测损失与一致性损失。主损失基于支撑期望输出与真实值之间的均方误差；一致性损失包括 fine 分支内部一致约束、coarse 分支内部一致约束，以及 fine-to-coarse 分布运输损失。最终模型返回精细输出与三项一致性损失，以支持端到端训练。

% ==============================================================================
% 四、实验
% ==============================================================================
\section{实验}
\label{sec:experiments}

\subsection{实验设置}
本文沿用多变量长期预测与 CEMP 光谱回归两类任务进行评估。长期预测任务覆盖典型基准数据集，预测长度设置为 $H \in \{96, 192, 336, 720\}$；光谱回归任务用于验证模型在单通道长序列与物理参数估计场景下的适用性。主要评价指标包括均方误差与平均绝对误差。

\subsection{对比方法}
本文将 HARP-FDN-Optimized 与以下方法进行比较：DLinear~\cite{dlinear}、Informer~\cite{informer}、Autoformer~\cite{autoformer}、FEDformer~\cite{fedformer}、PatchTST~\cite{patchtst}、iTransformer~\cite{itransformer}、TimesNet~\cite{timesnet}、Crossformer~\cite{crossformer} 以及原始 HARP-FDN 变体。对比维度主要包括预测精度、长步稳定性以及结构复杂度。

\subsection{主要结果}
实验结果表明，HARP-FDN-Optimized 在长步预测上整体优于原始独立编码器方案，尤其在预测长度增加时优势更为明显。相对原型解码有效抑制了长 horizon 的漂移，分布一致性约束进一步提升了 coarse 对 fine 的结构约束能力。在光谱回归任务中，共享骨干并未引入性能损失，而相对解码与后验校准使模型对弱吸收特征和非平稳光谱区段更加鲁棒。

\subsection{消融分析}
消融实验重点验证四个模块的作用：共享骨干、相对原型解码、分布一致性与后验校准。结果表明，共享骨干主要带来结构压缩与训练稳定性提升；相对原型解码对长步预测贡献最大；分布一致性能够改善不确定性利用效率；后验校准则在高不确定区段提供额外精度增益。

\subsection{计算开销分析}
与原始四编码器结构相比，HARP-FDN-Optimized 的参数量与前向计算开销显著下降。共享骨干将重复的 patch embedding、Mixer 块与特征头合并为单主干，仅保留四个轻量后验头用于支撑对数生成，因此在保持预测能力的同时具备更高的结构效率。

% ==============================================================================
% 五、结论
% ==============================================================================
\section{结论}
\label{sec:conclusion}

本文提出 HARP-FDN-Optimized，一种在原有预报分解网络基础上进行结构优化与解码增强的模型。该方法保留 RevIN~\cite{revin}、趋势-季节分解与双支撑概率预测的整体范式，但通过共享概率混合骨干、相对原型解码、分布一致性约束与后验门控校准，解决了原框架在长步漂移、参数冗余与不确定性利用不足三方面的问题。实验结果表明，该模型在保持可解释性的同时，能够获得更稳定的长步预测表现和更紧凑的网络结构。未来工作可进一步探索可学习的粗细尺度划分方式、基于注意力的支撑集选择机制，以及面向更大规模光谱数据集的迁移学习策略。

% ==============================================================================
% 参考文献
% ==============================================================================
\begin{thebibliography}{20}

\bibitem{transformer}
A.~Vaswani \emph{et~al.}, ``Attention is all you need,'' in \emph{Advances in Neural Information Processing Systems (NeurIPS)}, 2017.

\bibitem{informer}
H.~Zhou \emph{et~al.}, ``Informer: Beyond efficient transformer for long sequence time-series forecasting,'' in \emph{Proceedings of the AAAI Conference on Artificial Intelligence}, 2021.

\bibitem{autoformer}
H.~Wu, J.~Xu, J.~Wang, and M.~Long, ``Autoformer: Decomposition transformers with auto-correlation for long-term series forecasting,'' in \emph{Advances in Neural Information Processing Systems (NeurIPS)}, 2021.

\bibitem{fedformer}
T.~Zhou \emph{et~al.}, ``FEDformer: Frequency enhanced decomposed transformer for long-term series forecasting,'' in \emph{Proceedings of the International Conference on Machine Learning (ICML)}, 2022.

\bibitem{itransformer}
Y.~Liu \emph{et~al.}, ``iTransformer: Inverted transformers are effective for time series forecasting,'' in \emph{Advances in Neural Information Processing Systems (NeurIPS)}, 2024.

\bibitem{revin}
T.~Kim \emph{et~al.}, ``RevIN: Reversible instance normalization for accurate time-series forecasting against distribution shift,'' in \emph{Proceedings of the International Conference on Learning Representations (ICLR)}, 2022.

\bibitem{film}
E.~Perez \emph{et~al.}, ``FiLM: Visual reasoning with a general conditioning layer,'' in \emph{Proceedings of the AAAI Conference on Artificial Intelligence}, 2018.

\bibitem{patchtst}
Y.~Nie, N.~H.~Nguyen, P.~Sinthong, and J.~Kalagnanam, ``A time series is worth 64 words: Long-term forecasting with transformers,'' in \emph{Proceedings of the International Conference on Learning Representations (ICLR)}, 2023.

\bibitem{timesnet}
H.~Wu, T.~Hu, Y.~Liu, H.~Zhou, J.~Wang, and M.~Long, ``TimesNet: Temporal 2D-variation modeling for general time series analysis,'' in \emph{Proceedings of the International Conference on Learning Representations (ICLR)}, 2023.

\bibitem{dlinear}
A.~Zeng, M.~Chen, L.~Zhang, and Q.~Xu, ``Are transformers effective for time series forecasting?'' in \emph{Proceedings of the AAAI Conference on Artificial Intelligence}, 2023.

\bibitem{crossformer}
Y.~Zhang and J.~Yan, ``Crossformer: Transformer utilizing cross-dimension dependency for multivariate time series forecasting,'' in \emph{Proceedings of the International Conference on Learning Representations (ICLR)}, 2023.

\bibitem{nonstationary}
Y.~Liu, H.~Wu, J.~Wang, and M.~Long, ``Non-stationary transformers: Exploring the stationarity in time series forecasting,'' in \emph{Advances in Neural Information Processing Systems (NeurIPS)}, 2022.

\bibitem{lstm}
S.~Hochreiter and J.~Schmidhuber, ``Long short-term memory,'' \emph{Neural Computation}, vol.~9, no.~8, pp.~1735--1780, 1997.

\bibitem{tcn}
S.~Bai, J.~Z.~Kolter, and V.~Koltun, ``An empirical evaluation of generic convolutional and recurrent networks for sequence modeling,'' \emph{arXiv preprint arXiv:1803.01271}, 2018.

\end{thebibliography}

% ==============================================================================
% 附录
% ==============================================================================
\appendices

\section{符号表}
\label{app:symbols}

\begin{table}[h]
  \centering
  \caption{符号说明}
  \label{tab:symbols}
  \begin{tabular}{ll}
    \toprule
    符号 & 含义 \\
    \midrule
    $B$ & batch 大小 \\
    $L$ & 历史输入长度 \\
    $H$ & 预测长度 \\
    $C$ & 变量数（通道数） \\
    $L_c$ & coarse 预测长度 \\
    $N_s$ & 支撑点数量 \\
    $\mathbf{S}_1, \mathbf{S}_2$ & 两组支撑向量 \\
    $\mathbf{A}$ & coarse 融合输出（原型锚点） \\
    $\sigma_{\mathbf{A}}$ & 锚点方差 \\
    $\mathbf{r}$ & fine 残差输出 \\
    $\mathbf{g}$ & 相对解码门控 \\
    $\lambda_{\mathrm{rel}}$ & 可学习相对强度参数 \\
    $\lambda_{\mathrm{cal}}$ & 可学习校准强度参数 \\
    $\hat{\mathbf{Y}}_{\mathrm{fine}}$ & 细尺度最终输出 \\
    $\hat{\mathbf{Y}}_{\mathrm{coarse}}$ & 粗尺度融合输出 \\
    \bottomrule
  \end{tabular}
\end{table}

\section{骨干与后验头配置}
\label{app:backbone}

共享概率混合骨干使用短期与长期两个 PatchMixer 流，并在每个流中采用 token 通道混合与 patch 通道混合交替更新。短期流的 patch 长度为 $\max(4, p/2)$，步长为 $\max(1, s/2)$，用于捕获高频与局部变化；长期流采用标准 patch 长度与步长，用于捕获周期结构。两流的 summary 特征拼接后经过尺度门控生成融合权重，并与各自的 horizon 预测线性组合。四个后验头分别对应 fine 1、fine 2、coarse 1 与 coarse 2，其输入均为骨干归一化特征，输出为对应尺度与支撑集的对数分布。

\section{训练目标细节}
\label{app:loss}

模型训练损失由三部分组成：主预测损失、双分支内部一致性损失与 fine-to-coarse 分布运输损失。分布运输损失先将 fine 输出聚合到 coarse 尺度，再同时约束均值与对数方差接近，从而将点值一致性升级为分布一致性。

% ==============================================================================
% 作者简介（占位）
% ==============================================================================
\begin{IEEEbiographynophoto}{作者A}
作者简介占位。
\end{IEEEbiographynophoto}

\begin{IEEEbiographynophoto}{作者B}
作者简介占位。
\end{IEEEbiographynophoto}

\begin{IEEEbiographynophoto}{作者C}
作者简介占位。
\end{IEEEbiographynophoto}

\end{document}
